%! Dividing Polynomials; Remainder and Factor Theorems — Unit 4, Lesson 4.3 — Slides (Beamer)
\documentclass[aspectratio=169, 11pt]{beamer}
\usepackage{algebra2-beamer}   % provides \forestheader{} and \sectionlabel[color]{}

\begin{document}

% TITLE SLIDE (CourseName is not defined in beamer, so the name is literal)
{
\setbeamercolor{background canvas}{bg=forest}
\begin{frame}[plain]
  \vspace{2.8cm}\hspace{0.55cm}%
  \begin{minipage}{0.9\textwidth}
    {\color{goldacc}\bfseries\small LESSON 4.3}\\[0.5cm]
    {\color{white}\bfseries\LARGE Dividing Polynomials}\\[0.25cm]
    {\color{white}\bfseries\large Remainder \& Factor Theorems}\\[1.0cm]
    {\color{forestmid}\small Algebra 2: Shepherd~$\cdot$~Unit 4: Polynomial Functions}
  \end{minipage}
\end{frame}
}

% HOOK
\begin{frame}[t, plain]
  \forestheader{The wall we hit yesterday}
  \sectionlabel{Hook}
  \begin{columns}[T]
    \begin{column}{0.52\textwidth}
      \[ p(x) = x^3 - 3x + 2 \]
      \vspace{2pt}
      \begin{itemize}\setlength{\itemsep}{4pt}
        \item GCF? No. \; Two-term pattern? No.
        \item Quadratic form? No. \; Four terms? No.
        \item And yet Lesson 4.0's graph \emph{touches} at $x=1$ and \emph{crosses} at $x=-2$.
      \end{itemize}
    \end{column}
    \begin{column}{0.46\textwidth}
      \begin{block}{From the Warm-Up}
        $p(1) = 0$ \quad and \quad $(x-1)(x^2+x-2) = x^3-3x+2$. \\[4pt]
        So $(x-1)$ \textbf{is} a factor.
      \end{block}
      \vspace{3pt}
      {\footnotesize But nobody handed you $x^2+x-2$. If you know one factor, how do you find the
       other? \textbf{You divide.}}
    \end{column}
  \end{columns}
\end{frame}

% THE DIVISION STATEMENT
\begin{frame}[t, plain]
  \forestheader{One statement runs the whole lesson}
  \sectionlabel[goldacc]{Foundation}
  \[ 247 = 5\cdot 49 + 2 \qquad\longrightarrow\qquad P(x) = D(x)\,Q(x) + R(x) \]
  \begin{itemize}\setlength{\itemsep}{5pt}
    \item Equivalently: \; $\dfrac{P(x)}{D(x)} = Q(x) + \dfrac{R(x)}{D(x)}$
    \item Stop dividing when $\deg R < \deg D$.
    \item Check \emph{any} division: multiply $D\cdot Q$, add $R$, compare.
  \end{itemize}
  \begin{block}{The line the rest of the unit hangs on}
    If $R = 0$, then $P(x) = D(x)\,Q(x)$ --- so $D(x)$ is a \textbf{factor} of $P(x)$.
  \end{block}
\end{frame}

% MONOMIAL DIVISORS
\begin{frame}[t, plain]
  \forestheader{Monomial divisors: split the fraction}
  \sectionlabel{Method}
  \begin{columns}[T]
    \begin{column}{0.52\textwidth}
      \footnotesize
      \[ \frac{a+b}{c} = \frac{a}{c} + \frac{b}{c} \]
      \[ \frac{12x^4-8x^3+6x^2}{2x^2} = 6x^2-4x+3 \]
      \[ \frac{15x^3y^2-10x^2y^3+5x^2y^2}{5x^2y^2} = 3x-2y+1 \]
    \end{column}
    \begin{column}{0.44\textwidth}
      \begin{block}{The trap}
        The last term cancels completely --- and leaves $\mathbf{1}$, not $0$. \\[4pt]
        A quantity divided by itself is $1$.
      \end{block}
    \end{column}
  \end{columns}
\end{frame}

% LONG DIVISION
\begin{frame}[t, plain]
  \forestheader{Long division: two setup rules, then the wall falls}
  \sectionlabel{Method}
  \begin{columns}[T]
    \begin{column}{0.5\textwidth}
      \footnotesize
      \textbf{1.} Standard form. \quad \textbf{2.} A placeholder $0$ for every missing power.
      \[
      \begin{array}{r}
        x^2 + x - 2 \phantom{{}+2} \\[2pt]
        x-1\,\overline{\smash{\big)}\,x^3+0x^2-3x+2} \\[2pt]
        \underline{-(x^3-\phantom{1}x^2)}\phantom{{}-3x+2} \\[2pt]
        x^2-3x\phantom{{}+2} \\[2pt]
        \underline{-(x^2-\phantom{1}x)}\phantom{{}+2} \\[2pt]
        -2x+2 \\[2pt]
        \underline{-(-2x+2)} \\[2pt]
        0
      \end{array}
      \]
    \end{column}
    \begin{column}{0.48\textwidth}
      \begin{block}{Remainder $0$ --- so cash it in}
        \[ x^3-3x+2 = (x-1)(x^2+x-2) \]
        \[ = (x-1)^2(x+2) \]
      \end{block}
      \vspace{2pt}
      {\footnotesize Multiplicity $2$ at $x=1$ $\Rightarrow$ the graph \emph{touches}; multiplicity
       $1$ at $x=-2$ $\Rightarrow$ it \emph{crosses}. Exactly the Lesson 4.0 graph.
       \textbf{The wall is down.}}
    \end{column}
  \end{columns}
\end{frame}

% SYNTHETIC DIVISION
\begin{frame}[t, plain]
  \forestheader{Synthetic division: the same work, variables erased}
  \sectionlabel[goldacc]{Shortcut}
  \begin{columns}[T]
    \begin{column}{0.52\textwidth}
      \footnotesize
      $x^3-3x+2 \;\div\; (x-1)$, coefficients $1,\,0,\,-3,\,2$:
      \begin{center}
      \renewcommand{\arraystretch}{1.4}
      \begin{tabular}{r|rrrr}
      $1$ & $1$ & $0$  & $-3$ & $2$ \\
          &     & $1$  & $1$  & $-2$ \\ \cline{2-5}
          & $1$ & $1$  & $-2$ & $0$ \\
      \end{tabular}
      \end{center}
      \vspace{3pt}
      Bring down, multiply by $r$, add. The last entry is the \textbf{remainder}; the rest are the
      quotient's coefficients, one degree lower.
    \end{column}
    \begin{column}{0.46\textwidth}
      \begin{block}{One restriction, one trap}
        Divisors of the form $x-r$ \textbf{only} --- that is why long division never retires. \\[4pt]
        $x+2$ means $r = \mathbf{-2}$. Rewrite it as $x-(-2)$ every single time.
      \end{block}
      \vspace{2pt}
      {\footnotesize And the placeholder $0$ is still required in the coefficient row.}
    \end{column}
  \end{columns}
\end{frame}

% REMAINDER THEOREM
\begin{frame}[t, plain]
  \forestheader{The remainder was never leftovers}
  \sectionlabel[goldacc]{Remainder Theorem}
  \begin{columns}[T]
    \begin{column}{0.5\textwidth}
      \footnotesize
      $2x^3+3x^2-5 \;\div\; (x+2)$ gives
      \[ 2x^2-x+2 \;-\; \frac{9}{x+2} \]
      Now just \emph{evaluate}:
      \[ f(-2) = 2(-8)+3(4)-5 = -9 \]
      \textbf{The same number.}
    \end{column}
    \begin{column}{0.48\textwidth}
      \begin{block}{Remainder Theorem}
        Dividing $f(x)$ by $x-r$ leaves the remainder $f(r)$.
      \end{block}
      \vspace{3pt}
      {\footnotesize \textbf{Why:} $f(x) = (x-r)q(x) + R$. Put $x=r$: \\[3pt]
       $f(r) = 0\cdot q(r) + R = R$. \\[5pt]
       It runs both ways --- a remainder without dividing, \emph{and} a fast way to evaluate $f$.}
    \end{column}
  \end{columns}
\end{frame}

% FACTOR THEOREM
\begin{frame}[t, plain]
  \forestheader{Four names for one fact}
  \sectionlabel[goldacc]{Factor Theorem}
  \begin{center}
  \renewcommand{\arraystretch}{1.5}
  \begin{tabular}{|c|l|}
  \hline
  1 & $f(r) = 0$ \\ \hline
  2 & The remainder on dividing $f(x)$ by $x-r$ is $0$ \\ \hline
  3 & $(x-r)$ is a \textbf{factor} of $f(x)$ \\ \hline
  4 & $r$ is a \textbf{zero} of $f$, and $(r,0)$ is an $x$-intercept of its graph \\ \hline
  \end{tabular}
  \end{center}
  \vspace{4pt}
  \begin{block}{The workflow for the rest of Unit 4}
    \textbf{Test} a candidate $r$ $\;\rightarrow\;$ \textbf{divide} it out $\;\rightarrow\;$ factor
    the \textbf{depressed polynomial} completely.
  \end{block}
\end{frame}

% THE NEXT WALL
\begin{frame}[t, plain]
  \forestheader{Where does the candidate list come from?}
  \sectionlabel{Connections}
  \begin{columns}[T]
    \begin{column}{0.5\textwidth}
      \footnotesize
      We only tried $r=1$ because the Warm-Up pointed us there. \\[6pt]
      Try it cold: \; $h(x) = 2x^3-3x^2-11x+6$. \\[4pt]
      Integer divisors of $6$: $\pm1,\pm2,\pm3,\pm6$ --- and $r=-2$, $r=3$ both work. \\[4pt]
      \[ h(x) = (x+2)(2x-1)(x-3) \]
      The third zero is $\tfrac12$. \textbf{It was never on the list.}
    \end{column}
    \begin{column}{0.48\textwidth}
      \textbf{Today:} divide by monomials, binomials, and trinomials; synthetic division for $x-r$;
      the Remainder \& Factor Theorems.\\[8pt]
      \begin{itemize}\setlength{\itemsep}{3pt}
        \item \textbf{4.4} The Rational Root Theorem --- the candidate list, fractions included.
        \item \textbf{4.5} Fundamental Theorem of Algebra \& complex zeros.
        \item \textbf{4.6} Graphing polynomial functions \& modeling.
      \end{itemize}
    \end{column}
  \end{columns}
\end{frame}

\end{document}
